\documentclass[11pt,a4paper]{article}

\usepackage[T1]{fontenc}
\usepackage[utf8]{inputenc}
\usepackage{lmodern}
\usepackage{microtype}
\usepackage[margin=24mm]{geometry}
\usepackage{amsmath,amssymb,amsthm,mathtools,bm}
\usepackage{booktabs,longtable,array,tabularx,multirow}
\usepackage{graphicx}
\usepackage{caption}
\usepackage{subcaption}
\usepackage{enumitem}
\usepackage{xcolor}
\usepackage{listings}
\usepackage{fancyhdr}
\usepackage{titlesec}
\usepackage{hyperref}
\usepackage[nameinlink,capitalise,noabbrev]{cleveref}
\usepackage[numbers,sort&compress]{natbib}
\usepackage{doi}

\hypersetup{
  colorlinks=true,
  linkcolor=blue!55!black,
  citecolor=blue!55!black,
  urlcolor=blue!55!black,
  pdftitle={Dirac, BRST, Kaluza-Klein and Global Audit of a Chern-Weil Parent for Einstein Gravity},
  pdfauthor={Angus Muffatti}
}

\pagestyle{fancy}
\fancyhf{}
\lhead{TP-01 v1.1}
\rhead{Dirac/BRST and global audit}
\cfoot{\thepage}
\setlength{\headheight}{14pt}

\setlist[itemize]{leftmargin=*,topsep=3pt,itemsep=2pt}
\setlist[enumerate]{leftmargin=*,topsep=3pt,itemsep=2pt}
\setlength{\parskip}{4pt}
\setlength{\parindent}{0pt}

\lstdefinestyle{tpcode}{
  basicstyle=\ttfamily\small,
  frame=single,
  breaklines=true,
  columns=fullflexible,
  showstringspaces=false
}
\lstset{style=tpcode}

\newtheorem{theorem}{Theorem}
\newtheorem{lemma}[theorem]{Lemma}
\newtheorem{proposition}[theorem]{Proposition}
\newtheorem{corollary}[theorem]{Corollary}
\theoremstyle{definition}
\newtheorem{definition}[theorem]{Definition}
\theoremstyle{remark}
\newtheorem{remark}[theorem]{Remark}

\newcommand{\dd}{\mathrm d}
\newcommand{\CS}{\operatorname{CS}}
\newcommand{\AdS}{\operatorname{AdS}}
\newcommand{\Spin}{\operatorname{Spin}}
\newcommand{\Tr}{\operatorname{Tr}}
\newcommand{\Lie}[1]{\mathfrak{#1}}
\newcommand{\dof}{N_{\mathrm{dof}}}
\newcommand{\cA}{\mathcal A}
\newcommand{\cF}{\mathcal F}
\newcommand{\Om}{\Omega}
\newcommand{\phiz}{\phi}
\newcommand{\LieD}{\mathcal L}
\newcommand{\eps}{\epsilon}

\newcommand{\statusbox}[2]{%
\begin{center}
\fcolorbox{#1}{#1!5}{\parbox{0.93\textwidth}{#2}}
\end{center}}

\title{\textbf{From Exact Descent to Failed Reduction}\\
\large Dirac, BRST, Kaluza--Klein and global audit of a relative Chern--Weil parent for Einstein gravity}
\author{Angus Muffatti}
\date{TP-01 version 1.1 --- 28 August 2026}

\begin{document}
\maketitle

\begin{abstract}
Version 1.0 constructed an exact action-level genealogy linking a six-dimensional relative Chern--Weil functional, a five-dimensional transgression, genuine five-dimensional anti-de Sitter Chern--Simons gravity, and a four-dimensional MacDowell--Mansouri sector that expands to Einstein--Cartan gravity with cosmological and Euler terms.  It left open one decisive question: whether the fixed-holonomy Einstein sector is an admissible dynamical reduction of the parent rather than a deletion of fields and equations.  We answer that question by carrying out the canonical Dirac analysis in a regular sector, constructing the corresponding BRST/BFV complex, retaining the full circle zero mode and the first nonzero Kaluza--Klein level, pulling back the covariant presymplectic potential, and auditing bundles, Wilson lines and large gauge transformations.

The positive result is sharper than expected.  The pullback of the full zero-mode presymplectic potential to the fixed-holonomy surface has exactly the Einstein--Cartan plus Euler coefficients.  The failure is therefore not a missing gravitational kinetic coefficient.  The negative result is decisive.  The circle holonomy has gauge-invariant conjugacy data; its fixed norm is not a BRST gauge condition.  Varying it retains a gauge-invariant curvature-square equation that generic Einstein solutions violate.  A regular five-dimensional $\Spin(4,2)$ Chern--Simons sector carries thirteen local configuration degrees of freedom.  The maximally symmetric Chern--Simons vacuum is a rank-changing singular stratum with a vanishing ordinary quadratic graviton operator.  The nonlinear mode set $n=0,\pm1$ is not BRST closed, and its first conjugate Kaluza--Klein pair already adds twenty-six real regular-sector degrees of freedom at quadratic order.

We strengthen the restricted no-go theorem accordingly: under local, polynomial, pure-connection, regular, all-zero-mode, BRST-preserving and globally declared compactification assumptions, no such parent reduces to pure four-dimensional Einstein gravity with only two graviton helicities.  Exact cohomological genealogy survives; healthy spectrum-preserving dynamical equivalence does not.
\end{abstract}

\statusbox{blue}{\textbf{Terminal result.}
\[
\boxed{
\begin{array}{c}
\text{Exact common geometric genealogy: YES.}\\[1mm]
\text{Correct Einstein-sector symplectic pullback: YES.}\\[1mm]
\text{Dynamically invariant BRST/symplectic reduction: NO.}\\[1mm]
\text{Healthy two-helicity parent in the stated class: NO.}
\end{array}}
\]
The named v1.0 calculation is complete.  No high-performance computation is required for this verdict.}

\tableofcontents
\newpage

\section{Provenance, question and scope}

The motivating interview proposed that Chern--Simons/transgression and Einstein--Hilbert actions might be ``daughters'' of a deeper theory.  The interview supplied a question, not a mathematical construction or evidence.  TP-01 therefore separated two logically different claims:

\begin{description}[leftmargin=3.2cm,style=nextline]
\item[Weak genealogy claim] A single geometric object admits controlled descent, specialization, dimensional reduction or explicitly labelled sector restriction to both action types.
\item[Strong dynamical claim] A regular parent has a phase space, constraint algebra and physical spectrum that reduce to both daughters without deleting equations, imposing external gauge-invariant data or hiding modes at a singular locus.
\end{description}

Version 1.0 established the weak claim and rejected the strong claim provisionally using the extra Wilson-line equation, the generic degree count and the degenerate Chern--Simons vacuum.  It named the precise remaining calculation: retain all zero modes and the first nonzero Kaluza--Klein level, derive the Dirac and BRST systems, test preservation of the fixed-holonomy surface, compare presymplectic forms, and audit global bundles and large gauge transformations.  This paper performs that calculation.

The TP-00 acceptance framework is imported without re-interpretation.  In particular, a fatal algebraic or perturbative failure is not averaged against exact differential-form identities or attractive phenomenology.  The primary relevant gates are the theory tuple, symmetry closure, global structure, physical degree count, constraint algebra, perturbative health, radiative/EFT closure and Einstein recovery.

\paragraph{Scope limit.}
The canonical result is complete for the claimed reduction and for a generic regular canonical sector of five-dimensional Chern--Simons theory.  It is not an exhaustive classification of every degenerate and irregular phase-space stratum.  That limitation weakens no conclusion: the proposed Einstein vacuum lies at a known singular stratum, while a regular stratum has thirteen modes rather than two.

\section{Conventions and complete theory tuple}

All manifolds are smooth and oriented.  Four-dimensional spacetime has signature $(-,+,+,+)$ and $c=\hbar=1$.  Wedge products are suppressed where unambiguous.  Internal indices $\widehat A,\widehat B=0,\ldots,5$ are raised with
\[
\eta_{\widehat A\widehat B}=\operatorname{diag}(-,+,+,+,+,-),
\qquad
\epsilon_{012345}=+1.
\]
The gauge algebra is $\Lie{so}(4,2)$, with generators
\[
[J_{\widehat A\widehat B},J_{\widehat C\widehat D}]
=\eta_{\widehat B\widehat C}J_{\widehat A\widehat D}
-\eta_{\widehat A\widehat C}J_{\widehat B\widehat D}
-\eta_{\widehat B\widehat D}J_{\widehat A\widehat C}
+\eta_{\widehat A\widehat D}J_{\widehat B\widehat C}.
\]
